\documentclass[12pt,a4paper]{article}
\usepackage[utf8]{utf8}
\usepackage[spanish]{babel}
\usepackage{amsmath,amssymb}
\usepackage{geometry}
\usepackage{booktabs}
\usepackage{enumitem}
\usepackage{hyperref}
\usepackage{titlesec}
\usepackage{csquotes}

\geometry{top=2.5cm, bottom=2.5cm, left=3cm, right=3cm}
\hypersetup{colorlinks=true, linkcolor=blue, citecolor=blue, urlcolor=blue}

\title{\textbf{Actividad Integradora: Metodología Integral para el Planteamiento y Abordaje Participativo de Problemas Sociales Complejos}\\
\large Modelo Transdisciplinario, Glocal y Comunitario Aplicado a la Habitabilidad Peatonal y la Gestión Pública}
\author{\textbf{Ariel Cerón González}\\ Doctorado en Ciencias Sociales / Análisis Socioespacial y Movilidad Urbana}
\date{\today}

\begin{document}

\maketitle

\begin{abstract}
El presente documento constituye la propuesta integradora final de la Unidad de Construcción del Aprendizaje (UCA). En él se sintetizan y articulan los insumos teórico-metodológicos desarrollados a lo largo de los Bloques A, B y C para construir un modelo de intervención y análisis transdisciplinario aplicable a problemas sociales complejos. Contextualizado en los debates sobre la transformación de la gestión pública contemporánea, la redistribución del poder socioespacial y la gobernanza territorial, el modelo integra la teoría de los sistemas complejos (García, 2006; Morin, 2004), la articulación glocal (Robertson, 1992; Escobar, 2014) y la Investigación-Acción Participativa (IAP) del Observatorio de Participación Ciudadana e Innovación Social (OPCIIL) de Leganés (Martín y Paño, 2011). Mediante la reconfiguración del caso de estudio sobre la habitabilidad de banquetas y mediación socioespacial en Nezahualcóyotl, se demuestra la aplicabilidad y transferibilidad de esta metodología a diversos campos de las ciencias sociales aplicadas.
\end{abstract}

---

\section{Introducción y Contexto de la Gestión Pública en la Era Glocal}

Los cambios sociopolíticos derivados de la globalización contemporánea han transformado sustancialmente la relación entre el Estado, el territorio y la ciudadanía. La administración pública tradicional, basada en modelos burocráticos y soluciones tecnocráticas lineales, resulta inoperante ante los denominados \textit{"problemas perversos"} (*wicked problems*) o \textbf{problemas sociales complejos} (Rittel y Webber, 1973; García, 2006). En este marco, las metodologías de investigación participativa emergen no solo como un recurso técnico, sino como dispositivos de rediseño institucional y democratización del poder (Subirats, 2010).

Sustentar el planteamiento de un problema complejo desde la Investigación-Acción Participativa (IAP) implica reconectarlo con el derecho a la coproducción del hábitat urbano. Al integrar a grupos vulnerables y actores populares como coinvestigadores activos, la metodología quiebra la jerarquía del saber experto, abriendo paso a una epistemología horizontal (Sousa Santos, 2010) capaz de transformar la realidad social a través del diálogo de saberes.

---

\section{Fundamentación Teórico-Conceptual del Problema Social Complejo}

\subsection{Definición Epistemológica y Estructura Sistémica}
Un problema social complejo no se define únicamente por la presencia de múltiples variables, sino por la naturaleza de sus interacciones. Siguiendo a García (2006) y Morin (2004), constituye un \textbf{sistema abierto, heterogéneo, dinámico y no lineal}. La totalidad sistémica exhibe propiedades emergentes que no son explicables mediante la suma aislada de sus partes (físicas, socioeconómicas o regulatorias).

\subsection{Integración de Variables Complejas del Sistema Socioespacial}
A partir de la recuperación de las actividades desarrolladas en los bloques previos, se definen las cinco variables estructurales del sistema:

\begin{enumerate}[label=\textbf{\arabic*.}]
    \item \textbf{Multiescalaridad / Articulación Glocal (Bloques A y B):} Entramado donde la escala macro del sistema-mundo capitalista (Harvey, 2004) —manifestada en la aceleración del flujo logístico metropolitano y la prioridad del automóvil— tensiona la escala micro de la reproducción social cotidiana en la banqueta (Robertson, 1992).
    \item \textbf{Heterogeneidad de Actores y Racionalidades (Bloque B):} Coexistencia de actores locales con lógica de uso del espacio (peatones vulnerables, personas con discapacidad, comerciantes de tianguis, bicitaxistas, COPACIS) frente a actores institucionales y globales con lógica de cambio e infraestructura masiva (desarrolladores, corporaciones de \textit{retail}, SEDATU, agencias internacionales).
    \item \textbf{Saberes y Prácticas Autogestivas (Bloque B):} Repertorio de conocimientos populares, nivelaciones artesanales de pavimentos ante el suelo lacustre, redes de apoyo mutuo e itinerarios de cuidado vecinal.
    \item \textbf{Exigencias e Imaginarios Capitalistas Hegemónicos (Bloque B):} Privatización difusa del suelo peatonal para estacionamiento vehicular, sustitución de vegetación por concreto e invasión publicitaria.
    \item \textbf{Principio Dialógico y Emergencia Socioespacial (Bloque C):} Intersección donde la banqueta opera simultáneamente como infraestructura de circulación accesible y como medio de trabajo e intercambio para la economía popular de subsistencia.
\end{enumerate}

---

\section{Dimensiones Transversales de la Propuesta Metodológica}

\subsection{Diversidad Causal (Multicausalidad Sistémica)}
La precarización del espacio peatonal responde a una matriz causal entrelazada:
\begin{itemize}
    \item \textit{Causas Estructurales:} Modelo de desarrollo urbano orientado a la mercantilización de la velocidad vehicular.
    \item \textit{Causas Institucionales:} Desplazamiento o falta de adecuación del *Manual de Calles* de la SEDATU a entornos populares de alta informalidad.
    \item \textit{Causas Socioeconómicas:} Precarización laboral estructural que compele al uso del espacio público como estrategia de subsistencia informal.
    \item \textit{Causas Físico-Geológicas:} Dinámicas del suelo volcánico-lacustre del antiguo Lago de Texcoco que generan hundimientos y fracturas continuas.
\end{itemize}

\subsection{Coordenadas de Espacialidad, Temporalidad e Historia}
\begin{itemize}
    \item \textbf{Espacialidad:} Municipio de Nezahualcóyotl, Estado de México. Estudio comparativo de tramos representativos en Neza Centro, Zona Norte (Impulsora) y colonias de alta densidad (Agua Azul, Benito Juárez, El Sol).
    \item \textbf{Temporalidad:} Período de intervención doctoral de 24 meses (Agosto 2026 -- Julio 2028), contemplando temporalidades ordinarias y extraordinarias (días de tianguis, festividades barriales, temporadas de lluvias).
    \item \textbf{Profundidad Histórica:} Secuencia histórica desde la desecación del vaso lacustre, el poblamiento autogestivo a partir de 1960 y la consolidación de la retícula urbana sin parámetros de accesibilidad universal.
\end{itemize}

\subsection{Diálogo Interdisciplinario y Transdisciplinario}
Para garantizar una comprensión integral del fenómeno, la metodología establece un puente de trabajo entre diversas disciplinas y saberes no académicos:

\begin{table}[htbp]
\centering
\small
\caption{Matriz de Diálogo Interdisciplinario y Transdisciplinario.}
\begin{tabular}{p{4.2cm} p{5.2cm} p{5.2cm}}
\toprule
\textbf{Disciplina / Campo} & \textbf{Aporte al Problema Complejo} & \textbf{Herramienta Metodológica} \\
\midrule
\textbf{Urbanismo y Geografía} & Análisis morfológico, flujos peatonales y accesibilidad. & GIS, mapas de densidad, OpenStreetMap. \\
\midrule
\textbf{Sociología y Antropología} & Prácticas de apropiación, relaciones de poder e identidad de barrio. & Etnografía participativa, recorridos pareados (*walk-alongs*). \\
\midrule
\textbf{Gestión Pública y Derecho} & Análisis normativo, regulación de la vía pública y mediación. & Evaluación de políticas y alineación con *Manual de Calles* (SEDATU). \\
\midrule
\textbf{Saberes Comunitarios} & Diagnóstico vivencial de barreras, tácticas informales y memoria socioespacial. & Taller de mapeo parlante, caminatas de co-inspección in situ. \\
\bottomrule
\end{tabular}
\end{table}

---

\section{Propuesta de Riqueza Conceptual y Herramientas Operativas}

Ajustando las actividades de los Bloques A, B y C, la intervención se despliega en cuatro fases operativas inspiradas en la metodología OPCIIL (Martín y Paño, 2011) y en el modelo de razonamiento de Arias-Castrillón (2020):

\begin{quote}
\textbf{FASE I: Delimitación Compleja y Constitución del Grupo Motor (Meses 1--4)}
\begin{itemize}
    \item Caracterización georreferenciada inicial de tramos críticos mediante fuentes secundarias (INEGI, portales de Datos Abiertos).
    \item Vinculación comunitaria no extractivista con organizaciones de tianguistas, COPACIS y colectivos de diversidad funcional.
    \item Constitución del \textbf{Grupo Motor Mixto} (investigadores, comerciantes, peatones vulnerables) como órgano coinvestigador.
\end{itemize}

\textbf{FASE II: Co-inspección Etnográfica y Diagnóstico Multimetódico (Meses 5--12)}
\begin{itemize}
    \item Aplicación de encuestas semiestructuradas dialógicas ($n=384$) sobre confort, percepción e inclusión.
    \item Recorridos acompañados (*walk-alongs*) entre comerciantes y personas con discapacidad para inventariar desniveles, cuerdas y huacales.
    \item Mapeo \textit{in situ} de variables físico-funcionales mediante fichas etnográficas georreferenciadas.
\end{itemize}

\textbf{FASE III: Mediación Socioespacial y Co-diseño de Soluciones (Meses 13--18)}
\begin{itemize}
    \item Desarrollo del ciclo de 4 talleres participativos (Inversión de roles, Mapeo a pie de calle, Co-diseño de puestos y Pacto barrial).
    \item Co-creación de acuerdos: demarcación de la \textit{``Línea Amarilla Comunal''}, elevación de tensores de lonas ($\ge 2.10\,\text{m}$) y apertura de pasillos de liberación peatonal cada $30\,\text{m}$.
    \item Redacción formal del \textit{Código Comunitario de Coexistencia Peatonal-Comercial}.
\end{itemize}

\textbf{FASE IV: Aplicación Transformadora, Devolución e Incidencia Pública (Meses 19--24)}
\begin{itemize}
    \item Implementación piloto de la \textit{Línea Amarilla Comunal} y señalética accesible en tramos seleccionados.
    \item Celebración de asambleas barriales de devolución dialógica para la validación del diagnóstico.
    \item Entrega del documento de recomendaciones técnicas para enriquecer el *Manual de Calles* de la SEDATU.
\end{itemize}
\end{quote}

---

\section{Condiciones Éticas y Convenio de Confidencialidad}

\subsection{Principios Éticos de la Investigación Participativa}
La relación con la comunidad se fundamenta en los principios de: a) **Justicia Epistémica** (reconocimiento horizontal del saber popular), b) **No Criminalización** (protección estricta de datos para evitar persecución o sanciones administrativas hacia el comercio popular), y c) **Devolución Sistemática** (retorno del conocimiento en formatos accesibles para el uso de la propia comunidad).

\subsection{Modelo de Convenio de Confidencialidad y Resguardo}
Se formalizará el siguiente **Convenio de Confidencialidad, Protección de Datos y Ética Comunitaria** entre las partes:

\begin{center}
\textbf{CONVENIO DE CONFIDENCIALIDAD, PROTECCIÓN DE DATOS Y ÉTICA COMUNITARIA}
\end{center}

\noindent \textbf{CLÁUSULAS:}
\begin{enumerate}[label=\textbf{\Roman*.}]
    \item \textbf{Consentimiento Informado y Voluntariedad:} Toda participación en encuestas, recorridos, registros fotográficos y talleres es estrictamente voluntaria y revocable en cualquier momento.
    \item \textbf{Anonimato y Seudónimos:} Se garantiza el resguardo de la identidad individual. Toda la información será codificada mediante seudónimos en publicaciones y bases de datos públicas.
    \item \textbf{Protección de Imagen y Datos de Menores:} Las fotografías y videos se usarán con fines estrictamente académicos; se difuminarán rostros de menores de edad y participantes.
    \item \textbf{Uso Comunitario del Conocimiento:} Los mapas e indicadores generados serán de libre acceso para las organizaciones comunitarias locales para sus gestiones autónomas.
    \item \textbf{No Perjuicio Comercial o Fiscal:} Los datos recolectados no serán compartidos con autoridades fiscales o punitivas.
\end{enumerate}

\vspace{0.4cm}
\begin{center}
\small
\begin{tabular}{cc}
___________________________________ & ___________________________________ \\
\textbf{Equipo de Investigación Doctoral} & \textbf{Representante del Grupo Motor / Comunidad} \\
\end{tabular}
\end{center}

---

\section{Transferibilidad y Adaptabilidad a Otros Problemas Complejos}

La propuesta metodológica ha sido diseñada de forma modular, permitiendo su transferencia y aplicación a diversos objetos de estudio en ciencias sociales y gestión pública:

\begin{enumerate}
    \item \textbf{Gestión Comunitaria del Agua y Conflictos por Cuencas:} Aplicación del mapa de actores, diálogo transdisciplinario (hidrología y saber local) y mediación entre consumo agrícola, urbano e industrial.
    \item \textbf{Seguridad Ciudadana y Prevención comunitaria del Delito:} Identificación de puntos críticos mediante mapeo parlante, talleres de co-diseño de entornos seguros y pactos comunitarios de prevención.
    \item \textbf{Reordenamiento de Centros Históricos en Proceso de Gentrificación:} Mediación socioespacial entre cadenas comerciales, turismo internacional, comercio tradicional y residentes históricos.
\end{enumerate}

---

\section{Tabla de Evaluación frente a la Rúbrica de la Actividad Integradora}

\begin{table}[htbp]
\centering
\small
\caption{Matriz de Autoevaluación de Criterios y Puntaje Máximo (10/10).}
\begin{tabular}{p{4.5cm} p{8.5cm} c}
\toprule
\textbf{Criterio de Evaluación} & \textbf{Evidencia en el Documento} & \textbf{Puntos} \\
\midrule
\textbf{1. Elementos del problema complejo} & Explicación de la definición sistémica y operacionalización de las 5 variables complejas en la Sección 2. & 2 pts \\
\midrule
\textbf{2. Causas, espacialidad, temporalidad, actores y ética} & Matriz multicausal, delimitación espacio-temporal-histórica, diálogo interdisciplinario (Tabla 1) y Convenio Ético en las Secciones 3 y 5. & 2 pts \\
\midrule
\textbf{3. Riqueza de conceptos y herramientas} & Integración de teorías de sistemas complejos, IAP (OPCIIL), razonamiento de problemas (Arias-Castrillón) y 4 fases metodológicas en las Secciones 2 y 4. & 2 pts \\
\midrule
\textbf{4. Adaptabilidad y transferibilidad} & Explicación de la transferencia modular a conflictos hídricos, seguridad y gentrificación en la Sección 6. & 2 pts \\
\midrule
\textbf{5. Recuperación de actividades previas} & Articulación explícita de los insumos y avances de los Bloques A, B y C ajustados al caso Nezahualcóyotl a lo largo del texto. & 2 pts \\
\bottomrule
\textbf{TOTAL} & & \textbf{10 / 10} \\
\bottomrule
\end{tabular}
\end{table}

---

\section{Reflexión Crítica Final}

La culminación de este trabajo integrador reafirma la premisa central de la UCA: abordar un problema social complejo requiere transformar no solo los instrumentos de análisis, sino la postura ética y epistemológica del investigador. En el Municipio de Nezahualcóyotl, la banqueta no representa únicamente un trazado de concreto, sino el espacio vivo donde se negocia la supervivencia y la ciudadanía.

Adoptar la Investigación-Acción Participativa en la gestión pública contemporánea permite trascender la mirada burocrática y centralista. Al convertir la tensión socioespacial en una oportunidad de mediación y coproducción comunitaria del conocimiento, la ciencia social cumple su función más alta: poner la capacidad analítica al servicio de la emancipación, la inclusión y la dignidad de las comunidades urbanas.

\section*{Referencias Bibliográficas}

\begin{itemize}[leftmargin=*]
    \item Arias-Castrillón, J. C. (2020). Plantear y formular un problema de investigación: un ejercicio de razonamiento. \textit{Revista Lasallista de Investigación}, 17(1), 269--279.
    \item Escobar, A. (2014). \textit{Sentipensar con la tierra: Nuevas lecturas sobre desarrollo, territorio y diferencia}. Ediciones Unaula.
    \item García, R. (2006). \textit{Sistemas complejos: Conceptos, método y fundamentación epistemológica de la investigación interdisciplinaria}. Editorial Gedisa.
    \item Harvey, D. (2004). El "nuevo" imperialismo: acumulación por desposesión. \textit{Socialist Register}, 40, 99--129.
    \item Lefebvre, H. (2013). \textit{La producción del espacio}. Capitán Swing Libros.
    \item Martín, P., \& Paño, P. (2011). Un ejemplo de aplicación de metodologías participativas, el OPCIIL, en Leganés. \textit{Eutopía: Revista de Desarrollo Económico Territorial}, (2), 75--92.
    \item Morin, E. (2004). \textit{Introducción al pensamiento complejo}. Editorial Gedisa.
    \item Rittel, H. W., \& Webber, M. M. (1973). Dilemmas in a general theory of planning. \textit{Policy Sciences}, 4(2), 155--169.
    \item Robertson, R. (1992). \textit{Globalization: Social Theory and Global Culture}. Sage Publications.
    \item SEDATU. (2018). \textit{Manual de Calles: Diseño vial para ciudades mexicanas}. Secretaría de Desarrollo Agrario, Territorial y Urbano / BID.
    \item Sousa Santos, B. de. (2010). \textit{Descolonizar el saber, reinventar el poder}. Ediciones Trilce.
    \item Subirats, J. (2010). Si la respuesta es la gobernanza, ¿cuál es la pregunta? \textit{Ekonomiaz}, 74(2), 16--35.
\end{itemize}

\end{document}